\scanpage{116}
\uline{Lemme 1} $A$ anneau noethérien \deleted{intègre} réduit, $K$ son \deleted{corps} \inserted{anneau total} des fractions, $\widetilde A$ sa \deleted{\unsure{normalisation}} \inserted{clôture int.}

(i) [On suppose $\widetilde A=\displaystyle\bigcap_{p\text{ de rang }1}A_p$]. \deleted{\unsure{On suppose} $\widetilde A$ \unsure{entier sur} $A$}

Soit $f\in\deleted{\unsure{A}}\inserted{A\text{ non diviseur de zéro}}$, $p_i$ les idéaux pr. associés.

\deleted{(ii) Supposons $p_i A_{p_i}=fA_{p_i}$ (``multiplicités 1'')}
\margin{$fA=\bigcap p_i$}

(iii) $A/fA$ est réduit

\deleted{Soit $b\in A$}

Alors
\[
\boxed{A^*=A_f^*\cap\bigcap_i A_{p_i}^*.}
\]
\margin{\unsure{Sens des notations} si $A$ n'est pas intègre}

\uline{Dém} Posons $B=A_f\cap\displaystyle\bigcap_i A_{p_i}$. Soit $p$ pr. de rang 1; si $p\ni f$ on a $p=p_i$ pour $i$ \unsure{convenable}, d'où $A_p=A_{p_i}\supset B$; si $p\not\ni f$, alors $A_p\supset A_f\supset B$. Donc $B\subset\widetilde A$ par (i), \deleted{D'ailleurs $A_f^*\cap\displaystyle\bigcap_i A_{p_i}^*=B^*$.} Soit $x\in B$, donc $x=a/f^k$, $x$ entier sur $A$. On peut prendre $k\geq 0$ minimal. \deleted{Alors} On a
\[
x^n+c_1x^{n-1}+\cdots+c_n=0
\]
i.e.
\[
a^n+f^k c_1a^{n-1}+f^{2k}c_2a^{n-2}+\cdots+f^{nk}c_n=0
\]
d'où $a^n\in f^k A$; si $k\deleted{\unsure{\geq}} >0$, on aura donc $a^n\in fA=\bigcap p_i$, d'où $a\in\bigcap p_i\deleted{\unsure{\text{ d'où }a\in}}\inserted{=} fA$, \unsure{donc contraire au} \deleted{\unsure{dénominateur minimal}} \inserted{\unsure{caractère minimal de}} $a/f^k$
\margin{115}
